\begin{resumo}
\begin{SingleSpace}
Investidores utilizam índices de mercado como referência para análise de retorno, mas esses índices são generalistas, refletindo um comportamento médio. A construção de índices exige coleta, tratamento e composição de métricas sobre uma base de dados. Este trabalho propõe uma modelagem dimensional que representa dados de ações listadas na bolsa americana, e um ambiente para viabilizar a criação de novos índices que reflitam mais adequadamente os diferentes perfis de investidores. A avaliação foi feita através da comparação de carteiras compostas contra o índice S\&P 500, mostrando ganho de desempenho quando modificadores são incorporados ao critério de seleção. A principal contribuição deste trabalho é a integração entre coleta, modelagem e composição de métricas, e não apenas na geração de um indicador isolado. A curadoria do conjunto de ativos e a materialização das métricas derivadas na camada de dados asseguram a reprodutibilidade dos resultados e a auditabilidade do fluxo analítico.
\end{SingleSpace}
\vspace{\onelineskip}
\textbf{Palavras-chave}: análise de risco; mercado acionário; data mart; modelagem dimensional; smart beta.
\end{resumo}